\documentclass[../../../include/open-logic-section]{subfiles}

\begin{document}

\olfileid[pt]{sth}{story}{extensionality}
\olsection{Extensionalidade}

O primeiro ponto a notar é que os conjuntos são individuados pelos seus
!!{element}s. Mais precisamente:

\begin{axiom}[Extensionalidade]
Se os conjuntos $A$ e $B$ têm os mesmos !!{element}s, então $A$ e $B$ são
o mesmo conjunto.
\[
  \lforall[A][\lforall[B][(\lforall[x][(x \in A \liff x \in B)] \lif
  \eq[A][B])]]
\]
\end{axiom}

Assumimos isto ao longo de \olref[sfr][][]{part}. Mas vale a pena
repetir. O axioma de extensionalidade exprime a ideia fundamental de que um
conjunto é determinado pelos seus !!{element}s. (Assim, os conjuntos podem contrastar-se com os
\emph{conceitos}, nos quais precisamente os mesmos objetos podem cair sob
muitos conceitos diferentes.) 

Por que aceitar este princípio? Ora, é plausível dizer que qualquer
recusa da extensionalidade é uma decisão de abandonar tudo o que possa
ser chamado \emph{teoria dos conjuntos}. A teoria dos conjuntos não é mais nem menos do que
a teoria de coleções extensionais. 

O verdadeiro desafio na \olref[sth][][]{part}, contudo, é estabelecer
princípios que nos digam \emph{que conjuntos existem}. E acaba por se verificar que a única resposta verdadeiramente ``óbvia'' a esta questão é
provadamente errada.

\end{document}
